\documentclass[10pt]{article}
%\usepackage{epstopdf}
%\usepackage{fancyhdr}
\usepackage{graphicx}
\graphicspath{{./Figures/}}
\usepackage{fullpage}
\usepackage{bm}			% for bold greek letters
\hyphenation{}


\usepackage{tikz}
\usetikzlibrary{arrows,decorations.pathmorphing,backgrounds,positioning,fit}


\usepackage{amsmath}
\usepackage{amssymb}
\usepackage{amsthm}
\usepackage{array}
\usepackage{longtable}




\usepackage{subfigure}

%\linespread{1.3}
\author{Yang Li, Matthew Maron, Giacomo Po}
\title{Drift Diffusion}





\usepackage{hyperref}
\begin{document}
%%%%%%%%%%%%%%%%%%%%%%%%%%%%%%%%%%%%%%%%%%%%%
% TITLE
\maketitle
%%%%%%%%%%%%%%%%%%%%%%%%%%%%%%%%%%%%%%%%%%%%%

\section{Drift-Diffusion}

With $c$ expressed as a site fraction, define the normalized flux
$J_i=\Omega_0 j_i$, where $j_i$ is the defect-number flux and $\Omega_0$
is the atomic volume. The continuum normalized flux is
\begin{align}
J_i=-D_{ij}\left(c_{,j}+\frac{c}{kT}E^e_{,j}\right)
\end{align}
where $E^e$ is the interaction energy between the defect and the strain field at the equilibrium position.
Following, if we let $D_{ij}=\tilde{D}_{ij}\exp\left(\frac{E^e}{k_BT}\right)$ then we obtain:
\begin{align}
J_i=-\tilde{D}_{ij}\left(c_{,j}\exp\left(\frac{E^e}{k_BT}\right)+\frac{c}{kT}\exp\left(\frac{E^e}{k_BT}\right)E^e_{,j}\right)=-\tilde{D}_{ij}\left(c\exp\left(\frac{E^e}{k_BT}\right)\right)_{,j}=-\tilde{D}_{ij}\tilde{c}_{,j}
\end{align}
where we have defined the renormalized concentration
\begin{align}
\tilde{c}(\bm x)=c(\bm x)\exp\left(\frac{E^e(\bm x)}{k_BT}\right)
\end{align}
The renormalized diffusion tensor can be derived from  lattice theory \cite{Dederichs1978,DaFonseca2023}. It is
\begin{align}
\tilde{D}_{ij}(\bm x)=\frac{1}{4}D_0\sum_{\bm h} \hat{h}_i\hat{h}_j\exp\left(-\frac{E^{s,\bm h}(\bm x)}{k_BT}\right)
\label{Dt}
\end{align}
where $E^{s,\bm h}(\bm x)$ is the saddle-point interaction energy for a point defect initially located at $\bm x$ and performing a jump $\bm h$. 
Now consider the function $E$:
\begin{align}
E^{s,\bm h}(\bm x)=-P^{s,\bm h}_{ij}\varepsilon_{ij}(\bm x) = -P^{s,\bm h}_{ij}\left(\varepsilon^a_{ij}(\bm x) + \varepsilon^d_{ij}(\bm x)\right)
\end{align}
where $\bm P^{s,\bm h}$ is the elastic dipole tensor of the defect at the saddle point of the jump trajectory $\bm h$, and the elastic strain is decomposed in applied plus dislocation. The dipoe tensor is considered constant, that is its polarizability due to the srain field is neglected. 

First-order expansion of \eqref{Dt} yields
\begin{align}
\tilde{D}_{ij}(\bm x)=D_0\delta_{ij}+\underbrace{\frac{1}{4}\frac{D_0}{kT}\sum_{\bm h} \hat{h}_i\hat{h}_jP_{kl}^{s,\bm h}\varepsilon^a_{kl}(\bm x)}_{\text{SIPA-AD-Woo}}+\underbrace{\frac{1}{4}\frac{D_0}{kT}\sum_{\bm h} \hat{h}_i\hat{h}_jP_{kl}^{s,\bm h}\varepsilon^d_{kl}(\bm x)}_{\text{EID}}
\label{Dt0}
\end{align}

\paragraph{Diffusion far from dislocations}
Away from dislocations, and under a quasi uniform strain field, the renormalized diffusion tensor becomes a constant and generally anisotropic offset of the original diffusion coefficient
\begin{align}
\tilde{D}_{ij}=D_0\delta_{ij}+\frac{1}{4}\frac{D_0}{kT}\sum_{\bm h} \hat{h}_i\hat{h}_jP_{kl}^{s,\bm h}\varepsilon^a_{kl}
\end{align}
In this case, also the drift term in the original diffusion equation drops out, and the drift-diffusion problem reduces to a standard diffusoin problem with constant but stress-dependent diffusion coefficient. Anisotropic diffusion affects both diffusinal and dislocation creep. The effect on dislocation creep is due to the the fact that line elements receive a net flux that is different depending on their orientation in the anisotropic diffusion field. This is the Stress Induced Preferential Absorpition due to Anosotpropic Diffusion (SIPA-AD) proposed by Woo. Because the change in diffusion coefficient is proportional to the applied elastic strain, the stress exponent of SIPA-AD-Woo is expected to be 1.

SIPA-AD-Woo is a first order-effect in the defect dipole tensor. If the polarizability due to the strain field is considered, other second-order SIPA effects emerge (SIPA-I and SIPA-AD-Dederichs). 

 
\paragraph{Diffusion near dislocations}

\section{Modification of the diffusion--dislocation coupling}

For each mobile defect species $s$, let $\bm j_s$ denote the defect-number flux,
while $p_s$, $G_s$, and $R_s$ are rates of change of site fraction.  The steady
balance law is therefore written as
\begin{equation}
    -\Omega_0\bm\nabla\cdot\bm j_s
    +p_s+G_s+R_s(\bm c)=0,
    \label{eq:total_species_balance}
\end{equation}
where $p_s$ is the line-supported production term associated with dislocation
climb, $G_s$ is the irradiation production rate, and $R_s$ contains the bulk
first- and second-order reactions.  Introducing
\begin{equation}
    H_s(\bm x)
    =
    \exp\left(\frac{E_s^e(\bm x)}{k_BT}\right),
    \qquad
    W_s(\bm x)=H_s(\bm x)c_s(\bm x),
    \label{eq:activity_definition_theory}
\end{equation}
transforms the exact drift--diffusion flux into
\begin{equation}
    \bm j_s
    =
    -\frac{1}{\Omega_0}
    \widetilde{\bm D}_s(\bm x)\bm\nabla W_s,
    \qquad
    \widetilde{\bm D}_s(\bm x)
    =
    \widetilde{\bm D}_s
    \left(\bm\sigma_{\mathrm{ext}}
    +\bm\sigma_{\mathrm{DD}}(\bm x)\right).
    \label{eq:exact_activity_flux}
\end{equation}
The equilibrium-position interaction with the total stress is contained in
$H_s$, whereas $\widetilde{\bm D}_s$ describes the stress dependence of the
saddle-point migration process.  Equation~\eqref{eq:exact_activity_flux} still
contains the full diffusion tensor and no approximation has yet been made.

The activity is decomposed as
\begin{equation}
    W_s=W_s^\infty+\widetilde W_s,
    \label{eq:theory_activity_split}
\end{equation}
where $W_s^\infty$ is the dislocation-associated field and
$\widetilde W_s$ is a smooth finite-domain correction.  The physical total
concentration reconstructed from this decomposition remains
\begin{equation}
    c_s
    =
    H_s^{-1}
    \left(W_s^\infty+\widetilde W_s\right).
    \label{eq:theory_total_concentration}
\end{equation}
Because Eq.~\eqref{eq:exact_activity_flux} is linear in the activity gradient,
this decomposition induces the exact flux decomposition
\begin{align}
    \bm j_s
    &=
    \bm j_s^\infty+\widetilde{\bm j}_s,
    \label{eq:exact_flux_split}
    \\
    \bm j_s^\infty
    &=
    -\frac{1}{\Omega_0}
    \widetilde{\bm D}_s(\bm x)\bm\nabla W_s^\infty,
    \qquad
    \widetilde{\bm j}_s
    =
    -\frac{1}{\Omega_0}
    \widetilde{\bm D}_s(\bm x)\bm\nabla\widetilde W_s.
    \label{eq:exact_component_fluxes}
\end{align}
Both components therefore contain the same full tensor
$\widetilde{\bm D}_s(\bm\sigma_{\mathrm{ext}}+
\bm\sigma_{\mathrm{DD}})$ at this stage.  We define an exact source partition
by assigning the dislocation line source to the analytical component and the
smooth bulk terms to the finite-domain component:
\begin{align}
    -\Omega_0\bm\nabla\cdot\bm j_s^\infty+p_s&=0,
    \label{eq:exact_analytical_balance}
    \\
    -\Omega_0\bm\nabla\cdot\widetilde{\bm j}_s
    +G_s+R_s(\bm c)&=0.
    \label{eq:exact_smooth_balance}
\end{align}

The full tensor in Eq.~\eqref{eq:exact_component_fluxes} is strongly
heterogeneous near a dislocation and does not admit the constant-coefficient
anisotropic Green function used by the DDD formulation.  For a uniform applied
stress, let
$\bm D_s^{\mathrm{ext}}=\widetilde{\bm D}_s(\bm\sigma_{\mathrm{ext}})$
and approximate the analytical component flux as
\begin{equation}
    \bm j_s^{\infty,\mathrm{app}}
    =
    -\frac{Z_s^{\mathrm{EID}}}{\Omega_0}
    \bm D_s^{\mathrm{ext}}\bm\nabla W_s^\infty.
    \label{eq:biased_analytical_flux}
\end{equation}
Here $Z_s^{\mathrm{EID}}\bm D_s^{\mathrm{ext}}$ represents the residual effect
of the dislocation stress on saddle-point transport in the analytical channel;
the equilibrium-position interaction is already contained in $H_s$.
Substitution into the line-source balance gives
\begin{equation}
    \bm\nabla\cdot
    \left(
    \bm D_s^{\mathrm{ext}}\bm\nabla W_s^\infty
    \right)
    +
    \frac{p_s}{Z_s^{\mathrm{EID}}}
    =0.
    \label{eq:biased_analytical_problem}
\end{equation}
Thus the physical climb production $p_s$ remains conservative.  Let
$G_s^{\mathrm{ext}}$ denote the Green function with the convention
$\bm\nabla\cdot(\bm D_s^{\mathrm{ext}}\bm\nabla
G_s^{\mathrm{ext}})=-\delta$.  The analytical
solution can then be written as
\begin{equation}
    G_{Z,s}^{\mathrm{app}}
    =
    \frac{G_s^{\mathrm{ext}}}{Z_s^{\mathrm{EID}}},
    \qquad
    W_s^\infty
    =
    G_{Z,s}^{\mathrm{app}}*p_s.
    \label{eq:biased_green_approximation}
\end{equation}

For the smooth finite-domain component, the dislocation-scale variation of the
migration tensor is neglected and its exact flux is approximated as
\begin{equation}
    \widetilde{\bm j}_s^{\mathrm{app}}
    =
    -\frac{1}{\Omega_0}
    \bm D_s^{\mathrm{ext}}\bm\nabla\widetilde W_s.
    \label{eq:smooth_correction_flux}
\end{equation}
The exact reaction term in Eq.~\eqref{eq:exact_smooth_balance} depends on the
total concentration in Eq.~\eqref{eq:theory_total_concentration}.  To avoid
volume quadrature of the sharply peaked analytical field, it is separately
approximated by
\begin{equation}
    R_s(\bm c)
    \simeq
    R_s(\bm c^{\mathrm{bulk}}),
    \qquad
    c_s^{\mathrm{bulk}}
    =
    H_s^{-1}\widetilde W_s.
    \label{eq:bulk_reaction_concentration}
\end{equation}
The finite-domain problem is therefore
\begin{equation}
    \bm\nabla\cdot
    \left(
    \bm D_s^{\mathrm{ext}}\bm\nabla\widetilde W_s
    \right)
    +G_s
    +R_s(\bm c^{\mathrm{bulk}})
    =0
    \qquad\text{in }\Omega.
    \label{eq:smooth_fem_problem}
\end{equation}
The two component approximations define the reduced total flux
\begin{equation}
    \bm j_s^{\mathrm{red}}
    =
    \bm j_s^{\infty,\mathrm{app}}
    +\widetilde{\bm j}_s^{\mathrm{app}}
    =
    -\frac{\bm D_s^{\mathrm{ext}}}{\Omega_0}
    \left(
    Z_s^{\mathrm{EID}}\bm\nabla W_s^\infty
    +\bm\nabla\widetilde W_s
    \right),
    \label{eq:reduced_total_flux}
\end{equation}
and the closed reduced balance
\begin{equation}
    -\Omega_0\bm\nabla\cdot
    \left(
    \bm j_s^{\infty,\mathrm{app}}
    +\widetilde{\bm j}_s^{\mathrm{app}}
    \right)
    +p_s+G_s+R_s(\bm c^{\mathrm{bulk}})=0.
    \label{eq:reduced_species_balance}
\end{equation}
For $Z_s^{\mathrm{EID}}\ne1$, Eq.~\eqref{eq:reduced_total_flux} is not the
single local constitutive law
$-\bm D_s^{\mathrm{ext}}\bm\nabla(W_s^\infty+\widetilde W_s)/\Omega_0$.
Instead, it follows from applying different, explicitly stated approximations
to the two exact component fluxes in Eq.~\eqref{eq:exact_component_fluxes}.
Neglecting reactions involving $W_s^\infty$ assumes that diffusion and
absorption in the unresolved dislocation-core region are faster than the local
bulk reactions.  It also ensures that no strongly peaked infinite-domain
quantity enters the FEM volume quadrature.

If $W_s^{\mathrm b}$ denotes the prescribed boundary activity, the Dirichlet
condition for the correction is
\begin{equation}
    \left.\widetilde W_s\right|_{\partial\Omega}
    =
    W_s^{\mathrm b}
    -
    \left.W_s^\infty\right|_{\partial\Omega}.
    \label{eq:theory_dirichlet_split}
\end{equation}
The thermodynamic equilibrium condition on the dislocation line is retained
without modification:
\begin{equation}
    \left.
    \left(
    W_s^\infty+\widetilde W_s
    \right)
    \right|_{\mathcal N}
    =
    W_s^{\mathrm{eq}}.
    \label{eq:theory_dislocation_equilibrium}
\end{equation}
Thus the bias changes the climb flux associated with a prescribed activity
driving force, but it does not alter the equilibrium activity itself.

For a test function $\delta\widetilde W_s$ that vanishes on the Dirichlet
boundary, the weak form of Eq.~\eqref{eq:smooth_fem_problem} is
\begin{equation}
    \int_\Omega
    \bm\nabla(\delta\widetilde W_s)
    \cdot
    \left(
    \bm D_s^{\mathrm{ext}}\bm\nabla\widetilde W_s
    \right)
    \,\mathrm dV
    =
    \int_\Omega
    \delta\widetilde W_s
    \left[
    G_s+R_s(\bm c^{\mathrm{bulk}})
    \right]
    \,\mathrm dV.
    \label{eq:smooth_fem_weak_form}
\end{equation}
No analytical infinite-domain field appears in this volume integral.

\section{Dipole Tensor in Zr}

Point defects interact with the surrounding lattice through their elastic strain fields. At length scales larger than the defect core, this interaction can be represented by an elastic dipole tensor, $P_{ij}$, which describes the leading-order coupling between a point defect and an externally applied strain field: \begin{equation} 
\Delta E_{\mathrm{def}} = -P_{ij}\epsilon_{ij} 
\end{equation} 
$P_{ij}$ acts as a generalized force dipole.  Warwick, Ma, and Dudarev computed the elastic dipole tensors of hcp point defects using density functional theory. A point defect was introduced into a periodic supercell, and the atomic positions were relaxed while the cell shape and volume were held fixed. Because the defective crystal was not allowed to relax macroscopically, a residual average stress remained in the supercell. The elastic dipole tensor was then obtained from this residual stress as 
\begin{equation} 
P_{ij}(N_{\mathrm{def}}) = -V\overline{\sigma}_{ij}, 
\end{equation} 
where $V$ is the supercell volume and $\overline{\sigma}_{ij}$ is the average residual stress. The negative sign appears because the residual stress is the stress required to constrain the defective crystal at fixed cell shape and volume; the elastic dipole of the defect has the opposite sign. Since the calculation is performed in a finite periodic cell, the defect also interacts with its periodic images. To remove this finite-size effect, the dipole tensor components were computed for several supercell sizes and extrapolated to the isolated-defect limit using 
\begin{equation} 
P_{ij}(N) = \frac{\alpha_{ij}}{N} + P_{ij}^{\infty}, 
\end{equation} 
where $N$ is the number of atoms in the reference supercell, $\alpha_{ij}$ is a fitting constant, and $P_{ij}^{\infty}$ is the elastic dipole tensor of an isolated defect. The corresponding relaxation-volume tensor is related to the dipole tensor through 
\begin{equation} 
\Omega_{ij}=S_{ijkl}P_{kl}^{\infty}, 
\label{eq:relaxation_volume_tensor} 
\end{equation} 
where $S_{ijkl}$ is the elastic compliance tensor. Thus, $P_{ij}$ measures how the defect couples to strain or stress, while $\Omega_{ij}$ describes the macroscopic shape and volume change associated with the defect.

\subsection{Vacancy}
For a vacancy in Zr, the reported elastic dipole tensor is 
\begin{equation} 
\mathbf{P}^{\mathrm{Vac}} = 
\begin{pmatrix}
-5.156 & 0 & 0 \\ 
0 & -5.156 & 0 \\ 
0 & 0 & -8.022 
\end{pmatrix} \mathrm{eV}. 
\end{equation} 

All diagonal components are negative, indicating that the vacancy is a contractive defect. The contraction is anisotropic, and the vacancy produces a larger contractive distortion along the $c$-axis than in the basal plane. The corresponding relaxation-volume tensor is 
\begin{equation} 
\boldsymbol{\Omega}^{\mathrm{Vac}} = 
\begin{pmatrix} 
-0.085 & 0 & 0 \\
0 & -0.085 & 0 \\ 
0 & 0 & -0.267 
\end{pmatrix} \Omega_0
\end{equation} 
where $\Omega_0$ is the atomic volume of hcp Zr and $\Omega_{\mathrm{rel}}^{\mathrm{Vac}} = \mathrm{Tr}\left(\boldsymbol{\Omega}^{\mathrm{Vac}}\right) = -0.437\Omega_0$. Therefore, the vacancy reduces the crystal volume by approximately $0.44$ atomic volumes.

\subsection{Interstitial}

The lowest-energy self-interstitial configuration listed in the table is the basal-octahedral interstitial, denoted BO. Its elastic dipole tensor is 
\begin{equation} 
\mathbf{P}^{\mathrm{BO}} = 
\begin{pmatrix} 
16.655 & 0 & 0 \\ 
0 & 16.655 & 0 \\ 
0 & 0 & 10.344 
\end{pmatrix} \mathrm{eV}. 
\end{equation} 
All diagonal components are positive, showing that the basal-octahedral self-interstitial is an expansive defect. The expansion is anisotropic meaning that the defect couples more strongly to basal plane strain than to strain along the hcp $c$-axis. The corresponding relaxation-volume tensor is 
\begin{equation} 
\boldsymbol{\Omega}^{\mathrm{BO}} = 
\begin{pmatrix} 
0.549 & 0 & 0 \\ 
0 & 0.549 & 0 \\ 
0 & 0 & -0.008 
\end{pmatrix} \Omega_0,
\end{equation} 
where $\Omega_0$ is the atomic volume of hcp Zr. The scalar relaxation volume $\Omega_{\mathrm{rel}}^{\mathrm{BO}} = \mathrm{Tr}\left(\boldsymbol{\Omega}^{\mathrm{BO}}\right) = 1.09\Omega_0 $. Thus, the basal-octahedral self-interstitial increases the crystal volume by approximately one atomic volume. 

\section{Current Numerical Formulation and Consistency Audit}

This section records the formulation that is currently implemented in the
coupled cluster-dynamics and dislocation-dynamics code and identifies the
remaining limitations.  In particular, storing the finite-element unknown as
$W$ does not by itself transform a physical reaction operator; the required
mobile-reaction transformation is therefore included explicitly below.

\subsection{Nodal activity and the stress used in the transformation}

For each mobile species $s$, the code uses
\begin{equation}
    W_s = H_s c_s,
    \qquad
    H_s = \exp\left(\frac{E_s^e}{k_BT}\right).
    \label{eq:implemented_W_definition}
\end{equation}
At a finite-element node $A$, the stress entering $E_s^e$ is the patch-average
total stress
\begin{equation}
    \boldsymbol{\sigma}^{\mathrm{patch}}_A
    =
    \frac{1}{N_A}
    \sum_{e\in\mathcal{P}_A}
    \boldsymbol{\sigma}\left(\boldsymbol{x}_e^{c}\right),
    \label{eq:patch_average_stress}
\end{equation}
where $\mathcal{P}_A$ is the element patch of node $A$ and
$\boldsymbol{x}_e^{c}$ is the centroid of element $e$.  The stress returned by
the microstructure container is a sum of all pointwise stress contributions.
For the present simulations this is, to leading order,
\begin{equation}
    \boldsymbol{\sigma}^{\mathrm{patch}}_A
    =
    \boldsymbol{\sigma}^{\mathrm{DD,patch}}_A
    +
    \boldsymbol{\sigma}^{\mathrm{ext}}.
\end{equation}

The currently implemented equilibrium interaction energy is the hydrostatic,
scalar-relaxation-volume approximation
\begin{equation}
    E_{s,A}^e
    =
    -\frac{1}{3}
    \operatorname{tr}
    \left(\boldsymbol{\sigma}^{\mathrm{patch}}_A\right)
    n_s \overline{\Omega}^{\mathrm{rel}}_s,
    \qquad
    \overline{\Omega}^{\mathrm{rel}}_s
    =
    V^{\mathrm{rel}}_s\Omega_0,
    \label{eq:implemented_equilibrium_energy}
\end{equation}
where $n_s$ is the signed number of point defects in species $s$ and
$V^{\mathrm{rel}}_s$ is the input
\texttt{mobileSpeciesRelRelaxVol}.  Equation
\eqref{eq:implemented_equilibrium_energy} is internally consistent with the
current code, but it is not yet the full tensorial coupling
$E_s^e=-\boldsymbol{\sigma}:\boldsymbol{\Omega}_s$.  The anisotropic
relaxation-volume tensors discussed in the preceding section therefore remain
a future extension of the equilibrium interaction model.

Physical concentrations are converted to nodal activities when a
\texttt{cdMatrix} is read,
\begin{equation}
    W_{s,A}
    =
    c_{s,A}
    \exp\left(\frac{E_{s,A}^e}{k_BT}\right),
\end{equation}
and are converted back on output using the same patch-average stress,
\begin{equation}
    c_{s,A}
    =
    W_{s,A}
    \exp\left(-\frac{E_{s,A}^e}{k_BT}\right).
    \label{eq:nodal_output_conversion}
\end{equation}
Thus, input and output form an exact nodal round trip when the configuration
and stress state are unchanged.  Immobile-cluster variables are not affected
by this transformation.  During configuration loading, stress-producing
microstructures, including the current dislocation geometry, are initialized
before this nodal conversion is performed.  The conversion therefore uses the
stress of the frame being loaded rather than an empty or previous dislocation
configuration.

\subsection{External-stress-dependent diffusion tensor}

The diffusion tensor used by the finite-element transport operator depends
only on the macroscopic external stress.  Let
$\boldsymbol{Q}_g$ denote the crystal-to-global rotation of grain $g$.  The
stress is first rotated into the crystal frame,
\begin{equation}
    \boldsymbol{\sigma}^{g}_{\mathrm{ext}}
    =
    \boldsymbol{Q}_g^T
    \boldsymbol{\sigma}_{\mathrm{ext}}
    \boldsymbol{Q}_g.
    \label{eq:external_stress_local}
\end{equation}
For the component ordering
\begin{equation}
    \alpha\in\{11,12,13,22,23,33\},
\end{equation}
the implemented effective migration barrier and local diffusivity are
\begin{align}
    E^{m,\mathrm{eff}}_{s,\alpha}
    &=
    E^m_{s,\alpha}
    -
    \sigma^{g}_{\mathrm{ext},\alpha}
    V^{m}_{s,\alpha}\Omega_0,
    \label{eq:implemented_migration_barrier}
    \\
    \widetilde{D}^{g}_{s,\alpha}
    &=
    D^0_{s,\alpha}
    \exp\left(
    -\frac{E^{m,\mathrm{eff}}_{s,\alpha}}{k_BT}
    \right).
    \label{eq:implemented_local_diffusivity}
\end{align}
Here $V^m_{s,\alpha}$ is read from
\texttt{mobileSpeciesMigrationRelVolume}.  The global tensor is
\begin{equation}
    \widetilde{\boldsymbol{D}}^g_s
    =
    \boldsymbol{Q}_g
    \widetilde{\boldsymbol{D}}^{g,\mathrm{local}}_s
    \boldsymbol{Q}_g^T.
    \label{eq:implemented_global_diffusivity}
\end{equation}
No engineering-shear factor is applied to the stress components.  Therefore,
the migration-volume input must use the same component convention.  The
present implementation is a component-wise phenomenological
parameterization; a jump-resolved saddle-point model would instead evaluate a
scalar contraction for every jump direction and then perform the lattice sum
in Eq.~\eqref{Dt}.

The finite-element flux is
\begin{equation}
    \boldsymbol{j}_s^{\mathrm{FEM}}
    =
    -\frac{1}{\Omega_0}
    \widetilde{\boldsymbol{D}}^g_s
    \boldsymbol{\nabla}W_s^{\mathrm{CD}}.
    \label{eq:implemented_flux}
\end{equation}
Both the flux operator and its species scaling use
$\widetilde{\boldsymbol{D}}^g_s(\boldsymbol{\sigma}_{\mathrm{ext}})$.
When the external stress changes, the stored stress is updated and the FEM
matrix is marked for reassembly.  With the current uniform external-load
controller, the dislocation network, cluster dynamics, and inclusions have
zero average stress, so the stress supplied to the diffusion tensor is the
uniform external stress alone.  A spatially varying external elastic field is
not yet represented by this average-stress path.

\subsection{CD--DD superposition and boundary conditions}

The implemented decomposition is most accurately written as
\begin{equation}
    W_s
    =
    W_s^{\mathrm{CD}}
    +
    W_s^{\mathrm{DD}},
    \label{eq:implemented_superposition}
\end{equation}
where $W_s^{\mathrm{DD}}$ is the analytical dislocation contribution and
$W_s^{\mathrm{CD}}$ is the finite-element field.  The analytical contribution
uses an anisotropic Green-function kernel and includes the effective
dislocation bias through a multiplicative $1/Z_s^{\mathrm{EID}}$ factor.

The boundary condition is imposed directly as a Dirichlet condition in
activity space.  At a boundary node,
\begin{equation}
    W_{s}^{\mathrm{CD}}
    =
    W_{s}^{\mathrm{bnd}}
    -
    W_{s}^{\mathrm{DD}}.
    \label{eq:implemented_dirichlet_split}
\end{equation}
Consequently, the sum in Eq.~\eqref{eq:implemented_superposition} satisfies
the prescribed boundary activity.  The FEM correction is homogeneous only
when irradiation production and reactions are disabled.  In the general
case, the intended physical equation for the present split is
\begin{equation}
    \boldsymbol{\nabla}\cdot
    \left(
    \widetilde{\boldsymbol{D}}_s
    \boldsymbol{\nabla}W_s^{\mathrm{CD}}
    \right)
    +
    G_s
    +
    R_s(\boldsymbol{c}^{\mathrm{bulk}})
    =
    0
    \qquad
    \text{away from the analytical line sources}.
    \label{eq:implemented_cd_equation}
\end{equation}
The current code realizes Eq.~\eqref{eq:implemented_cd_equation} for both the
transport-only and reaction-enabled problems.  In the latter case, the smooth
physical concentration is reconstructed from the nodal activity before the
reaction operators are evaluated, as detailed below.

The analytical $W^{\mathrm{DD}}$ kernel uses the inverse and determinant of
the same external-stress-dependent tensor
$\widetilde{\boldsymbol{D}}(\boldsymbol{\sigma}_{\mathrm{ext}})$
as the FEM operator, together with the multiplicative
$1/Z_s^{\mathrm{EID}}$ Green-response factor.  These tensors are cached per
grain and mobile species.
The cache is rebuilt whenever the external stress changes, together with the
invalidation of the FEM solver matrix.  Thus, the analytical and
finite-element parts share the same external-stress-dependent diffusion
tensor, while the analytical branch additionally carries the effective bias
closure defined in Eq.~\eqref{eq:biased_green_approximation}.
The boundary degrees of freedom are registered once; later stress-triggered
matrix reassemblies reuse those constraints and only update their prescribed
activity values.

\subsection{Reaction terms in the activity formulation}

The finite-element unknown and its Newton increment are both stored as $W$.
The first- and second-order reaction operators remain physical concentration
operators; consequently, the code reconstructs the smooth bulk concentration
before evaluating either the reaction residual or its Newton linearization.

Define the diagonal nodal transformation
\begin{equation}
    \boldsymbol{H}_A
    =
    \operatorname{diag}
    \left[
    \exp\left(\frac{E_{1,A}^e}{k_BT}\right),
    \ldots,
    \exp\left(\frac{E_{m,A}^e}{k_BT}\right)
    \right].
\end{equation}
To remain consistent with the reduced model in
Eq.~\eqref{eq:smooth_fem_problem}, the smooth bulk concentration at a
quadrature point should be reconstructed from the nodal finite-element
activities,
\begin{equation}
    \boldsymbol{c}^{\mathrm{bulk}}_h(\boldsymbol{x})
    =
    \sum_A
    N_A(\boldsymbol{x})
    \boldsymbol{H}_A^{-1}
    \boldsymbol{W}^{\mathrm{CD}}_A.
    \label{eq:reaction_concentration_interpolation}
\end{equation}
For a physical reaction law
\begin{equation}
    R_a(\boldsymbol{c})
    =
    R^{(1)}_{ab}c_b
    +
    \frac{1}{2}
    R^{(2)}_{abg}c_bc_g,
    \label{eq:physical_reaction_law}
\end{equation}
the residual must evaluate Eq.~\eqref{eq:physical_reaction_law} using
Eq.~\eqref{eq:reaction_concentration_interpolation}.  Its derivative with
respect to nodal finite-element activity is
\begin{equation}
    \frac{\partial R_a}{\partial W^{\mathrm{CD}}_{s,B}}
    =
    \frac{\partial R_a}{\partial c_s}
    N_B
    \left(\boldsymbol{H}_B^{-1}\right)_{ss}.
    \label{eq:reaction_W_jacobian}
\end{equation}
The implementation applies the factors $\boldsymbol{H}_B^{-1}$ to the nodal
trial shape functions used by both the reaction residual and the reaction
Jacobian.  The quantity denoted by \texttt{R1cd} remains a separate
cluster-size scaling and is not used as the $c$-to-$W$ transformation.  The
strongly peaked analytical
$W^{\mathrm{DD}}$ field is deliberately excluded from this volume reaction
quadrature, consistently with the inner-reaction approximation in
Eq.~\eqref{eq:bulk_reaction_concentration}.

Reaction-rate coefficients also contain diffusivities.  The first- and
second-order coefficient matrices are cached per grain using
$\widetilde{\boldsymbol{D}}(\boldsymbol{\sigma}_{\mathrm{ext}})$, and the
special one-dimensional sink correction uses the corresponding crystal-frame
local diffusivity.  These caches are rebuilt together with the analytical
diffusion cache whenever the external stress changes.  The irradiation source
vector is also reassembled because its test-function scaling contains the same
stress-dependent diffusivity.

\subsection{Output and visualization}

The \texttt{cdMatrix} stores physical mobile concentrations.  DDqt follows
the same convention.  Before plotting, the nodal CD activity is converted
with Eq.~\eqref{eq:nodal_output_conversion}, and the resulting nodal physical
concentration is interpolated with the same quadratic finite-element shape
functions used by the solver.  The analytical Green response is returned
directly as $W_s^{\mathrm{DD}}$ and is converted at the visualization point
using the total point stress.  Therefore, the plotted quantities are
\begin{equation}
    c_s^{\mathrm{plot}}
    =
    c_s^{\mathrm{CD}}
    +
    c_s^{\mathrm{DD}},
\end{equation}
and the CD contribution agrees with the output \texttt{cdMatrix} at every
finite-element node.
Before a DDqt frame is initialized, its run ID, physical time, and time step are
restored from the corresponding row of the history file.  A generated frame
zero with no history records uses the initial time zero and initial step;
missing later frames remain errors.  Time-dependent
external loads and the resulting diffusion tensor are therefore evaluated at
the displayed frame time.

\subsection{Implementation status and verification sequence}

The present status can be summarized as follows:
\begin{center}
\renewcommand{\arraystretch}{1.2}
\begin{tabular}{|p{0.47\textwidth}|p{0.43\textwidth}|}
\hline
Component & Status \\
\hline
Nodal $c\leftrightarrow W$ conversion using patch-average total stress
& Implemented for input, equilibrium initialization, output, and CD visualization. \\
\hline
FEM transport with
$\widetilde{\boldsymbol{D}}(\boldsymbol{\sigma}_{\mathrm{ext}})$
& Implemented, including grain-frame rotation and matrix reassembly after an external-stress change. \\
\hline
Analytical DD kernel with the same external-stress-dependent tensor
& Implemented; its inverse and determinant tensors are refreshed together with the FEM operator after an external-stress change. \\
\hline
Diffusion-tensor admissibility
& Implemented with a finite-entry and positive-definiteness check before FEM or Green-kernel use. \\
\hline
First- and second-order reactions in the $W$ formulation
& Implemented using nodal physical-concentration reconstruction and the Jacobian factors in Eq.~\eqref{eq:reaction_W_jacobian}; the analytical DD field is excluded from bulk reaction quadrature. \\
\hline
External-stress-dependent mobile reaction coefficients
& Implemented per grain and refreshed after an external-stress change. \\
\hline
DDqt physical-concentration visualization
& Implemented consistently with nodal output, including displayed-frame time synchronization for time-dependent loads. \\
\hline
\end{tabular}
\end{center}

\subsection{Code changes for collaborator review}

Table~\ref{tab:code_changes} lists the functional code changes in the present
revision.  Generated build files and the removed historical
\texttt{GalerkinClimbSolverOld} sources are not included.

{\small
\begin{longtable}{|>{\raggedright\arraybackslash}p{0.28\textwidth}|>{\raggedright\arraybackslash}p{0.39\textwidth}|>{\raggedright\arraybackslash}p{0.23\textwidth}|}
\caption{Code changes implementing the activity formulation and
external-stress-dependent diffusion.}
\label{tab:code_changes}\\
\hline
Code location & Implemented change & Primary review point \\
\hline
\endfirsthead
\multicolumn{3}{c}{\tablename\ \thetable{} -- continued from previous page}\\
\hline
Code location & Implemented change & Primary review point \\
\hline
\endhead
\hline
\multicolumn{3}{r}{Continued on next page}\\
\endfoot
\hline
\endlastfoot

\texttt{ClusterDynamicsParameters.h}\newline
\texttt{ClusterDynamicsParameters.}\newline
\texttt{cpp}: parameter initialization and
\texttt{diffusionTensor()}
&
Added \texttt{mobileSpeciesMigrationRelVolume}; evaluated the
stress-dependent migration barriers in the crystal frame using component
order $[11,12,13,22,23,33]$; rotated the resulting diffusion tensor back to
the global frame; rejected non-finite or non-positive-definite tensors before
they reach the FEM and Green operators.
&
Check migration-volume units, barrier sign, component order, and
crystal-to-global rotation.
\\
\hline

\texttt{ClusterDynamicsParameters.h}\newline
\texttt{ClusterDynamicsParameters.}\newline
\texttt{cpp}: activity conversion helpers
&
Centralized $E^e$, $c\leftrightarrow W$, and equilibrium conversions.
\texttt{physicalMobileConcentration()} now evaluates the direct inverse
$W/\exp(E^e/k_BT)$ without an exponent clamp.
&
Check the sign of $E^e$ and the scalar relaxation-volume approximation.
\\
\hline

\texttt{ClusterDynamicsFEM.h}\newline
\texttt{ClusterDynamicsFEM.cpp}: transport operator and external stress
&
Made both \texttt{FluxMatrix} and \texttt{InvDscaling} use
$\widetilde{\bm D}(\bm\sigma_{\mathrm{ext}})$.  Added
\texttt{setExternalStress()}, FEM-matrix invalidation, and per-grain
\texttt{analyticalInvD}/\texttt{analyticalDetD} caches for the analytical
Green function.  The irradiation source vector is reassembled after an
external-stress change because it contains the same inverse-diffusivity test
scaling.
&
Check that only the macroscopic external stress enters diffusion and that a
stress change triggers cache refresh and matrix reassembly.
\\
\hline

\texttt{ClusterDynamicsFEM.h/.cpp}: Dirichlet registration
&
Separated one-time boundary-degree-of-freedom registration from repeat FEM
matrix assembly.  External-stress changes now rebuild the matrix without
reinserting zero conditions over the current nonzero activity values.
&
Check repeated external-stress updates with assertions enabled.
\\
\hline

\texttt{ClusterDynamicsParameters.h/.cpp}: reaction coefficients
&
Constructed first- and second-order mobile reaction coefficients from the
per-grain tensor $\widetilde{\bm D}(\bm\sigma_{\mathrm{ext}})$ rather than a
zero-stress diffusion cache.  The crystal-frame tensor used by the optional
one-dimensional sink correction is evaluated at the same external stress.
&
Check the scalar averaging of the anisotropic diffusivity and cache refresh
after a stress change.
\\
\hline

\texttt{ClusterDynamicsFEM.h/.cpp}\newline
\texttt{SecondOrderReaction.h}: mobile reactions
&
Reconstructed $\bm c_h^{\mathrm{bulk}}=\sum_A N_A\bm H_A^{-1}\bm W_A^{\mathrm{CD}}$
from patch-average total stress and used it in first- and second-order reaction
residuals.  Applied the same nodal factors to the Newton increment, thereby
including the chain-rule factor in the reaction Jacobian.
&
Check the nodal stress convention, exclusion of $W^{\mathrm{DD}}$, and finite-
difference agreement of the reaction Jacobian.
\\
\hline

\texttt{ClusterDynamicsFEM.cpp}: \texttt{initializeConfiguration()}
&
Interpreted input \texttt{cdMatrix} values as physical concentrations and
converted them to nodal activities.  Applied the same conversion to default
equilibrium initialization using the node patch-average total stress.
&
Check input dimensions and the element-centroid patch average at boundary and
interior nodes.
\\
\hline

\texttt{DefectiveCrystal.cpp}: configuration initialization
&
Initialized all stress-producing microstructures before Cluster Dynamics so
that input $c\rightarrow W$ conversion uses the current DD geometry rather
than an empty or previous-frame network.
&
Check first-load and DDqt frame-to-frame input/output round trips.
\\
\hline

\texttt{ClusterDynamics.h}\newline
\texttt{ClusterDynamics.cpp}: boundary conditions
&
Applied the Dirichlet split
$W^{\mathrm{CD}}=W^{\mathrm b}-W^{\mathrm{DD}}$ at boundary nodes.  The
prescribed physical boundary concentration is converted to $W^{\mathrm b}$
using the node patch-average total stress; the analytical
$W^{\mathrm{DD}}$ returned by the dislocation network is then subtracted.
&
Check the sign of the superposition, the prescribed boundary concentration,
and the point-stress versus patch-average-stress convention for
$W^{\mathrm{DD}}$ at boundary nodes.
\\
\hline

\texttt{ClusterDynamics.h}\newline
\texttt{ClusterDynamics.cpp}: solve and output paths
&
Updated the external stress before FEM assembly.  Added
\texttt{nodalPhysicalMobileConcentration()} and wrote physical mobile
concentrations, rather than activities, to \texttt{cdMatrix}.
&
Check CD-before-DD solve ordering and the nodal
$c\rightarrow W\rightarrow c$ round trip.
\\
\hline

\texttt{DislocationSegment.h}\newline
\texttt{DislocationSegment.cpp} and
\texttt{GalerkinClimbSolver.cpp}
&
Passed the runtime inverse and determinant of
$\widetilde{\bm D}(\bm\sigma_{\mathrm{ext}})$ into the anisotropic analytical
kernel and the Galerkin climb stiffness assembly.  Returned the Green response
directly as $W^{\mathrm{DD}}$, removing the former extra multiplication by
$H$.
&
Check that the analytical $W^{\mathrm{DD}}$ and FEM correction use the same
diffusion operator.
\\
\hline

\texttt{DislocationNetwork.cpp}
&
Used the runtime diffusion cache in the analytical mobile-activity field,
treated the Green response directly as $W^{\mathrm{DD}}$, and converted
quadrature-point CD/DD activities to physical concentrations using total
point stress on output.
&
Check the distinction between external stress in $\bm D$ and total stress in
$c\leftrightarrow W$.
\\
\hline

\texttt{DDFieldWidget.h/.cpp}\newline
\texttt{DefectiveCrystalActor.cpp}
&
Changed DDqt concentration fields from activity to physical concentration.
The CD contribution is reconstructed by interpolating nodal physical
concentrations with the quadratic FEM shape functions; other activity
contributions are converted at the visualization point using total stress.
The displayed frame time is restored from the history file before evaluating
external loads.
&
Check agreement with \texttt{cdMatrix} at FEM nodes and visualization-mesh
convergence away from nodes.
\\

\end{longtable}
}

The mobile reaction residual, reaction Jacobian, and diffusion-dependent
reaction coefficients are now consistent with the reduced activity
formulation.  Immobile-species evolution remains outside the present change.
The zero-stress \texttt{vClimbRef} is retained because it is used only as a
numerical time-step reference and does not enter the climb operator.

The recommended verification and development order is:
\begin{enumerate}
    \item Verify nodal $c\rightarrow W\rightarrow c$ round trips at fixed
    stress to machine precision.
    \item Verify that zero external stress reproduces the original diffusion
    tensors and that prescribed uniaxial and hydrostatic stresses reproduce
    Eqs.~\eqref{eq:implemented_migration_barrier}--\eqref{eq:implemented_global_diffusivity}.
    \item Check symmetry, positive definiteness, finiteness, and crystal-frame
    covariance of every diffusion tensor over the intended stress range.
    \item Confirm that changing the external stress invalidates and
    reassembles the FEM matrix.
    \item Verify the reaction residual and its analytical Jacobian against
    finite differences of Eqs.~\eqref{eq:reaction_concentration_interpolation}
    and \eqref{eq:reaction_W_jacobian}.
    \item Verify that all diffusion-dependent reaction coefficients are
    refreshed from the prescribed external stress.
    \item Compare DDqt nodal values against \texttt{cdMatrix} output and
    perform CD-mesh and visualization-mesh convergence studies.
    \item Finally, verify tensile/compressive symmetry, far-field behavior,
    positivity of physical concentration, and convergence of coupled
    dislocation-climb benchmark problems.
\end{enumerate}

\begin{thebibliography}{9}

\bibitem{Dederichs1978}
Dederichs, P. H., and K. Schroeder. "Anisotropic diffusion in stress fields." Physical Review B 17, no. 6 (1978): 2524.

  \bibitem{DaFonseca2023}
  Da Fonseca, Daphné, Fabien Onimus, Frédéric Mompiou, M-C. Marinica, Edouard de Sonis, Emmanuel Clouet, and Thomas Jourdan. "Numerical investigation of dislocation climb under stress and irradiation." Acta Materialia 242 (2023): 118431.
  
\end{thebibliography}

\end{document}
